\chapter{Objectives}
\label{chap:objectives}

The experiment has four objectives. They are listed in the order in which a
failure of one invalidates the ones below it, which is also the order in which
the session is run.

\section{O1. Establish a valid calibration on the integrated support}

The stereo geometry that every reconstructed number depends on has to be
measured on the support that will carry the cameras during the experiment, not
on any other mount. The archived calibration is unusable and the project knows
why: measured over 424 archived frames the median board tilt is
4.4\textdegree, only 39\% of frames exceed 5\textdegree, and not one frame in
the entire archive reaches 45\textdegree, which leaves the focal length
mathematically unrecoverable \cite{linku_calibration_method}. The guided tool
built in response refuses to report a session complete until the accepted
views cover four tilt bins.

This objective is a precondition, not a result. If it fails, objectives O2 to
O4 still produce images, but every metric derived from them is scaled by an
unknown factor and no distance in metres can be quoted.

\section{O2. Confirm the illumination reproduces the orbital condition}

One direction of light, no fill, no bounce, a parallel enough beam, and enough
of it. The requirements and their derivation are in the dark-room deliverable
\cite{linku_darkroom_illumination}; the measurement of whether the TV2 fixtures
actually meet them has never been made, because the room has never been
occupied. The enclosure used until now contains one unspecified bulb whose
irradiance, spectrum and beam behaviour have never been measured.

\section{O3. Run the full pipeline end to end on a real moving tether}

Segment, find the endpoints, trace the path in both images, rectify, match,
triangulate, and compute the metric scalars. Every stage has been exercised on
archived or simulated data. None has been exercised on a tether that was
moving, in a dark room, with the rig also moving.

Two specific failure modes are being looked for, both of which the project has
predicted but never observed in hardware:

\begin{itemize}
    \item \textbf{The parametric matching assumption.} Stage 5 pairs the point
    at $i\%$ along the left path with the point at $i\%$ along the right path.
    This holds only if both traced paths start and end at the same physical
    points on the tether and traverse it at the same rate. When an endpoint is
    occluded, or a path leaves the frame, the error is systematic and silent:
    there is no epipolar check to catch it. A hand-moved tether is the first
    situation in which an endpoint can plausibly leave the frame mid-sequence.

    \item \textbf{Rolling-shutter shear.} At the software's current 15 frame/s
    configuration the 3-pixel shear limit corresponds to roughly
    1.05\textdegree/s \cite{linku_rollingshutter}. Hand motion will exceed
    that. The experiment does not measure the rate, so it cannot quantify the
    shear; what it can establish is whether the tracker and the matching
    survive it at all, and at what point they visibly stop doing so.
\end{itemize}

\section{O4. Record the scene well enough to rebuild it in Blender}

The Blender validation deliverable compares a real photograph against a render
of the same scene, and it is currently blocked on two things this session can
supply: measured laboratory-condition light values, which stand in the scene as
placeholders, and the physical pose of the tether, which for the existing
reference pair was never written down
\cite{linku_blender_validation}. Chapter \ref{chap:record} is written to close
both.

\section{What this experiment deliberately does not do}
\label{sec:why-by-hand}

\textbf{It does not measure angular rate.} There is no servo, no motor and no
rotation stage. Both the tether and the camera rig are moved by hand,
erratically. A quantitative rolling-shutter characterisation needs a commanded
or encoded rate, and that is a separate experiment scheduled in a later phase
\cite{linku_rollingshutter}. Running it here with uninstrumented motion would
produce numbers that look quantitative and are not, which is worse than
producing none.

The choice of hand motion is not merely a concession to the absence of a rig.
The deployment it stands in for is itself not a smooth commanded rotation, and
the question this experiment asks is the coarser one that has to be answered
first: does the pipeline survive uncontrolled motion, or does it fail, and
where.

\textbf{It does not attempt a reconstruction accuracy figure against ground
truth.} No independently measured tether geometry is available at the scale
and in the pose used. What can be checked is internal consistency and physical
plausibility, which is what Chapter \ref{chap:criteria} registers.

\textbf{It does not use the flight sensor.} The IMX708 B0308 is the decided
flight sensor but has not been purchased, with procurement expected the
following year. This experiment uses the IMX477 pair that exists. Every
optical figure quoted in this document is therefore computed for the IMX477
with its 1.55~\textmu m pixel and the fitted 8~mm lens, and does not transfer
to the B0308.
